\documentclass[11pt,a4paper]{article}
\usepackage[T1]{fontenc}
\usepackage{lmodern}
\usepackage[margin=27mm]{geometry}
\usepackage{amsmath,amssymb,amsthm}
\usepackage{microtype}
\usepackage{booktabs}
\usepackage[numbers,sort&compress]{natbib}
\usepackage[hidelinks]{hyperref}
\usepackage{url}
\numberwithin{equation}{section}
\newtheorem{theorem}{Theorem}
\newcommand{\Rplus}{\mathbb{R}_{>0}}
\newcommand{\dd}{\mathrm{d}}
\pdfinfoomitdate=1
\pdftrailerid{}
\pdfsuppressptexinfo=-1
\hypersetup{pdftitle={Absence of reflection-symmetric two-periodic boundary laws in the competing three-state SOS model},pdfauthor={Xiaojian Zeng}}
\title{Absence of reflection-symmetric two-periodic boundary laws in the competing three-state SOS model}
\author{Xiaojian Zeng\\[3pt]
\small PBC School of Finance, Tsinghua University, 100190\\
\small \href{mailto:xiaojianzeng@foxmail.com}{xiaojianzeng@foxmail.com}}
\date{}

\begin{document}
\maketitle
\begin{abstract}
We consider the three-state solid-on-solid model on a rooted binary tree with
nearest-neighbor and sibling interactions, each sibling bond counted once.
For arbitrary finite real couplings and positive finite inverse temperature,
we prove that the exact recursion of positive reflection-symmetric boundary
messages has no nonfixed two-cycle. This resolves the two-periodicity
conjecture of Karshiboev and Rahmatullaev in its stated invariant subspace.
The proof reduces two distinct cycle values to their sum and product.
Positivity would force a negative linear coefficient, but in that entire
parameter region a strict inequality makes the corresponding discriminant
negative. We also give a square-plus-positive-remainder certificate for the
composition quotient. The boundary-message kernels yield compatible Gibbs
cylinder distributions, so the result excludes alternating generation-parity
boundary laws in this class. It does not assert uniqueness or exclude Gibbs
measures outside the reflection-symmetric, level-homogeneous class.
\end{abstract}

\noindent\textbf{Keywords:} solid-on-solid model; binary tree; competing interactions;
Gibbs measures; periodic boundary laws

\section{Introduction}
Periodic boundary messages provide a precise way to test for alternating
equilibrium states on a rooted tree. In the three-state solid-on-solid (SOS)
model with an additional interaction between siblings, the two children
can remain correlated after conditioning on their parent. Nevertheless,
summing over each parent--children unit gives an exact finite-dimensional
recursion. Its reflection-symmetric subspace reduces the question of
generation-parity alternation to a scalar two-cycle equation.

Karshiboev and Rahmatullaev derived this recursion and excluded nonfixed
two-cycles in a parameter region described by a nonnegative polynomial
coefficient $B$ \cite[Sec.~5, Theorem~3]{KR2025}. Their Conjecture~1 leaves
the complementary region $B<0$ to be excluded. Fixed points are omitted
from the two-periodic solutions under discussion, as specified immediately
after their Eq.~(16). The invariant set is $v=1$, where $v$ is the ratio
of the boundary weights for spins $2$ and $0$.

Here we close that remaining region for all positive values of both
Boltzmann parameters. The decisive observation is that $B<0$ forces
$t>1$ and $|B|<t^4-1$, whereas the leading and constant coefficients
obey $A>t^6$ and $C>t^8$. These inequalities contradict the existence
of two real positive cycle values. The proof is exact and requires no
parameter scan or limiting approximation. An alternative global
positivity certificate is given in Appendix~\ref{sec:certificate}.

\section{Model and boundary messages}
\label{sec:model}
Let every vertex of a rooted binary tree have two children, including
the root. Write $y_x,z_x$ for the children of $x$ and let
$\sigma_x\in\{0,1,2\}$. The interaction is
\begin{equation}
 \mathcal H(\sigma)
 =-J\sum_{\langle x,y\rangle}|\sigma_x-\sigma_y|
  -J_1\sum_x|\sigma_{y_x}-\sigma_{z_x}|.
 \label{eq:hamiltonian}
\end{equation}
The first sum counts every parent--child bond once; the second counts
the single unordered sibling pair of each parent once. Infinite-volume
Gibbs conditionals use the finite sum of interactions meeting the
conditioned volume. We assume finite $J,J_1\in\mathbb R$ and
$0<\beta<\infty$, and set
\begin{equation}
 t=e^{\beta J},\qquad h=e^{\beta J_1}.
 \label{eq:parameters}
\end{equation}
Thus $t,h\in\Rplus$ are independent. The couplings have energy units,
$\beta$ has inverse-energy units, and the spin differences and all message
ratios are dimensionless. No external field or other distance-two
interaction is included. These conventions agree with
\citet[Eq.~(1)]{KR2025}, with $(t,h)=(\theta,\theta_1)$.

The Boltzmann factor of a parent of spin $i$ and its labeled children
of spins $j,k$ is
\begin{equation}
 \psi(i,j,k)=t^{|i-j|+|i-k|}h^{|j-k|}.
 \label{eq:factor}
\end{equation}
For a positive child message $m$, define
\begin{equation}
 T_i(m)=\sum_{j,k=0}^{2}\psi(i,j,k)m(j)m(k).
 \label{eq:update}
\end{equation}
The two assignments $(j,k)$ and $(k,j)$ are distinct when $j\ne k$;
this does not double the sibling energy. Descendant subtrees communicate
only through the displayed local units, so \eqref{eq:update} is an exact
partition-weight recursion.

Normalize a reflection-symmetric message to $m=(1,u,1)$ with $u>0$.
Reflection of all three spins gives $T_2=T_0$. Direct evaluation of the
nine summands for each parent spin gives
\begin{align}
 N(u):=T_1(1,u,1)&=u^2+4thu+2t^2(h^2+1),\label{eq:N}\\
 D(u):=T_0(1,u,1)&=t^2u^2+2th(t^2+1)u+t^4+2t^2h^2+1.
 \label{eq:D}
\end{align}
Consequently the normalized update is
\begin{equation}
 f(u)=\frac{N(u)}{D(u)},\qquad u>0.
 \label{eq:f}
\end{equation}
Both numerator and denominator are strictly positive. This is the map
in \citet[Eq.~(12), p.~137]{KR2025}.

A positive level-homogeneous boundary law means a sequence of messages
$m_n$ shared by all vertices in generation $n$, defined up to positive
scalar multiples, with
\begin{equation}
 T(m_{n+1})=\lambda_n m_n,\qquad \lambda_n>0.
 \label{eq:compatibility}
\end{equation}
On the reflection-symmetric subspace this is $u_n=f(u_{n+1})$.
An alternating assignment $(1,a,1),(1,b,1)$ therefore requires
\begin{equation}
 a=f(b),\qquad b=f(a),\qquad a,b>0.
 \label{eq:cycle}
\end{equation}
A nonfixed two-cycle has $a\ne b$.

\section{Exclusion of two-cycles}
\label{sec:proof}
\begin{theorem}
\label{thm:main}
For every $t,h,u>0$, the function \eqref{eq:f} satisfies
\begin{equation}
 f(f(u))=u\quad\Longrightarrow\quad f(u)=u.
 \label{eq:theorem}
\end{equation}
Hence the model \eqref{eq:hamiltonian} has no nonconstant
generation-parity boundary law among positive reflection-symmetric
level-homogeneous messages.
\end{theorem}

\begin{proof}
Write
\begin{equation}
\begin{gathered}
 N(u)=u^2+n_1u+n_0,\qquad D(u)=t^2u^2+d_1u+d_0,\\
 n_1=4th,\quad n_0=2t^2(h^2+1),\quad L=d_0+n_1,\\
 d_1=2th(t^2+1),\qquad d_0=t^4+2t^2h^2+1.
\end{gathered}
\label{eq:coefficients}
\end{equation}
Assume that $a,b>0$ solve \eqref{eq:cycle} with $a\ne b$.
Set $F=aD(b)-N(b)$ and $G=bD(a)-N(a)$, so $F=G=0$.
Polynomial expansion yields the two identities
\begin{align}
 F-G&=(a-b)(L+a+b-t^2ab),\label{eq:antisym}\\
 F+G+(a+b)(L+a+b-t^2ab)
   &=2\bigl[d_0(a+b)+(d_1+1)ab-n_0\bigr].
 \label{eq:sym}
\end{align}
Division of \eqref{eq:antisym} by $a-b$ is legitimate under the
assumption of a nonfixed cycle. With $s=a+b$ and $p=ab$, the two
identities become
\begin{equation}
 t^2p=s+L,\qquad d_0s+(d_1+1)p=n_0.
 \label{eq:sumproductsystem}
\end{equation}
Define
\begin{equation}
 A=t^2d_0+d_1+1,\quad
 B=(d_1+1)L-t^2n_0,\quad
 C=d_0L+n_0.
 \label{eq:ABC}
\end{equation}
Solving \eqref{eq:sumproductsystem} gives
\begin{equation}
 s=-\frac{B}{A},\qquad p=\frac{C}{A}.
 \label{eq:sumproduct}
\end{equation}
Here $A>t^6$ since $d_0>t^4$, and $C>t^8$ since
$L>d_0>t^4$. Moreover,
\begin{equation}
\begin{split}
 B={}&1-t^4+2th(t^6+t^4+t^2+3)\\
    &+2t^2h^2(3t^2+5)+4t^3h^3(t^2+1)>1-t^4.
\end{split}
\label{eq:Bbound}
\end{equation}
Because $s>0$ and $A>0$, equation \eqref{eq:sumproduct} forces
$B<0$. Inequality \eqref{eq:Bbound} then implies $t>1$ and
$0<-B<t^4-1$. Therefore
\begin{equation}
 B^2<(t^4-1)^2<t^8<4t^{14}<4AC.
 \label{eq:chain}
\end{equation}
On the other hand, two real numbers $a,b$ with sum and product
\eqref{eq:sumproduct} necessarily satisfy
\begin{equation}
 A^2(a-b)^2=B^2-4AC\ge0,
 \label{eq:contradiction}
\end{equation}
contradicting \eqref{eq:chain}. No nonfixed positive two-cycle exists.
Finally, if $f(f(u))=u$, positivity of $f(u)$ allows
$a=u$, $b=f(u)$ in this argument. They must coincide.
\end{proof}

The coefficients \eqref{eq:ABC} define the quadratic obtained by
eliminating fixed points from the second iterate of \eqref{eq:f};
compare \citet[Eq.~(17), p.~139]{KR2025}. When $B\ge0$, $A,C>0$ already
exclude positive roots. Equation \eqref{eq:chain} excludes the entire
remaining region $B<0$, thus resolving Conjecture~1 in the nonfixed
sense of its preceding discussion. This includes the boundary $B=0$.
The statement is not that $B^2-4AC$ is negative for all positive $t,h$;
the sign implication is needed only when $B<0$.

\section{Gibbs interpretation and scope}
\label{sec:gibbs}
For completeness, the compatibility condition \eqref{eq:compatibility}
has a direct infinite-tree interpretation despite the sibling coupling.
For any compatible positive messages, put
\begin{equation}
 K_n(j,k\mid i)=
 \frac{\psi(i,j,k)m_{n+1}(j)m_{n+1}(k)}
      {\lambda_n m_n(i)}.
 \label{eq:kernel}
\end{equation}
Its sum over $j,k$ is one by \eqref{eq:compatibility}. Give the root
spin the law $m_0(i)/\sum_r m_0(r)$, and multiply these conditional
pair kernels over the parents in successive generations.

Let $W_n$ denote generation $n$ and $V_n=\bigcup_{k=0}^n W_k$.
The messages at internal vertices cancel, leaving the cylinder density
\begin{equation}
 \mu_n(\sigma_{V_n})=\frac{1}{Z_n}
 \prod_{x\in V_{n-1}}
 \psi(\sigma_x,\sigma_{y_x},\sigma_{z_x})
 \prod_{x\in W_n}m_n(\sigma_x),
 \label{eq:cylinder}
\end{equation}
where $Z_n>0$ normalizes the density. Summing a terminal generation
recovers $\mu_{n-1}$ because each kernel is normalized. These
consistent distributions on finite spin spaces define an infinite
probability measure.

To check the Gibbs condition for a finite vertex set $\Lambda$,
choose $V_n$ large enough that all local units meeting $\Lambda$
lie inside it and the terminal generation is outside $\Lambda$.
Conditioning \eqref{eq:cylinder} on the spins outside $\Lambda$
cancels terminal weights and all interactions disjoint from $\Lambda$.
The remaining factor product is exactly the Boltzmann weight of every
interaction meeting $\Lambda$, including a sibling bond crossing its
boundary. This conditional depends on only finitely many exterior
spins; the same identity persists in the consistent infinite measure.
Thus these boundary laws give the Gibbs specification of
\eqref{eq:hamiltonian}.

Theorem~\ref{thm:main} excludes an alternating assignment in this
positive reflection-symmetric class. It does not classify all Gibbs
measures or establish uniqueness, extremality, stability, or convergence
of arbitrary message iterations. In particular, reflection symmetry and
level homogeneity are hypotheses; the theorem does not exclude
two-periodic messages with unequal spin-$0$ and spin-$2$ weights.
This agrees with the scope distinction in
\citet[Remark~5(a)]{KR2025}. No zero-temperature conclusion is taken.

The nearest-neighbor limit also fixes the message convention. When
$h=1$,
\begin{equation}
 f(u)=\left(\frac{u+2t}{tu+t^2+1}\right)^2.
 \label{eq:nearest}
\end{equation}
The Hamiltonian of the later thermodynamic study
\cite[Eq.~(2.1), p.~166]{RK2026} has only nearest-neighbor interactions.
After shifting its spin labels by one, its branch recursion is
$g_{n+1}=W(g_n^2)$, where $W_{ij}=t^{|i-j|}$ and the square is
componentwise \cite[Eq.~(2.4), p.~166]{RK2026}.
Our descendant message satisfies $T(m)=(Wm)^2$ at $h=1$, so the
convention change is $m=g^2$. This one-coupling limit does not cover
independent $h\ne1$; no other thermodynamic or dynamical conclusions
of that study are used here.

\section{Exact verification}
\label{sec:verification}
The accompanying code contains two exact representations of the
calculation. The integer checker reconstructs all $27$ local weights
in \eqref{eq:factor}, forms two generations of the unnormalized
update, and verifies the composition identity in
Appendix~\ref{sec:certificate} by sparse polynomial arithmetic over
the integers. It uses no computer algebra library or cached
coefficient data. The symbolic checker reconstructs the local update
as a quadratic form with sibling matrix $S_{jk}=h^{|j-k|}$ and
independently derives \eqref{eq:antisym}--\eqref{eq:sumproduct}.
Both verify the positive remainder in \eqref{eq:positiveQ}.

As a finite-volume check of \eqref{eq:cylinder}, the symbolic-check
program also uses exact rational arithmetic at $t=3/2$, $h=2/5$ and
terminal message $(1,5/7,1)$. On a seven-vertex tree the kernel product
agrees with the direct normalized Hamiltonian weight for all $3^7=2187$
configurations. Its first-generation marginals and all $3^6=729$
exterior choices for an internal-site conditional retain the crossing
sibling interaction. These finite checks test the implementation; the
infinite-volume statement follows from the consistency argument above.
The algebraic identities hold for all parameters in their stated domain,
without floating-point error, truncation, or extrapolation.

\appendix
\section{A positive composition quotient}
\label{sec:certificate}
An alternative certificate avoids dividing by a fixed-point factor.
With the notation of \eqref{eq:coefficients} and \eqref{eq:ABC}, set
\begin{equation}
\begin{gathered}
 P(u)=N(u)-uD(u),\qquad Q(u)=Au^2+Bu+C,\\
 E(u)=t^2N(u)^2+d_1N(u)D(u)+d_0D(u)^2>0.
\end{gathered}
\label{eq:PQE}
\end{equation}
Direct polynomial multiplication gives
\begin{equation}
 N^2+n_1ND+n_0D^2-uE=PQ,
 \qquad f(f(u))-u=\frac{P(u)Q(u)}{E(u)}.
 \label{eq:composition}
\end{equation}
These are identities including at $P=0$. All coefficients are fixed
by the defining map; in particular the expansion of $C$ contains
$4t^6h^2$.

Separate the terms of positive degree in $h$:
\begin{align}
 Q(u)&=Q_0(u)+R_h(u),\label{eq:Qsplit}\\
 Q_0(u)&=(t^6+t^2+1)u^2+(1-t^4)u+t^8+2t^4+2t^2+1,
 \label{eq:Qzero}\\
 R_h(u)&=4t^4h^4+h^3(4t^5u+4t^3u+8t^3)\notag\\
 &\quad+h^2(2t^4u^2+6t^4u+10t^2u+4t^6+6t^2)\notag\\
 &\quad+h(2t^3u^2+2tu^2+2t^7u+2t^5u
              +2t^3u+6tu+4t^5+4t).
 \label{eq:Rh}
\end{align}
Every monomial in $R_h$ is positive for $t,h,u>0$. Completing the
single negative term in $Q_0$ into a square gives
\begin{equation}
\begin{split}
 Q(u)={}&t^2\left(t^2u-\tfrac12\right)^2
 +(t^2+1)u^2+u+t^8+2t^4+\tfrac74t^2+1+R_h(u)>1.
\end{split}
\label{eq:positiveQ}
\end{equation}
Here the formal substitution $h=0$ used to define $Q_0$ is only an
algebraic device; it does not introduce a zero Boltzmann weight into
the physical domain. Equations \eqref{eq:composition} and
\eqref{eq:positiveQ} imply $f(f(u))=u$ if and only if $P(u)=0$,
which gives a second proof of Theorem~\ref{thm:main}. The certificate
covers $t=1$, $B=0$, repeated fixed points, and any removable
algebraic factors of $f$ without cancellation at a fixed point.

\section*{Data availability}
All defining mathematical inputs are given in this article. The accompanying
reproducibility files contain the integer and symbolic checkers, their
environment requirements, and instructions for regenerating the exact
outputs. No empirical data or fitted parameters are used.

\section*{Acknowledgements and AI disclosure}
OpenAI Codex was used for proof exploration, algebraic reasoning, code
generation, source checking, internal computational review, and manuscript
drafting. This assistance included substantive work on the derivation.
The exact program checks described above do not constitute external peer
review. Human critical review and final approval of this draft are pending.

\bibliographystyle{unsrtnat}
\bibliography{references}
\end{document}
